% chapters/02_metodologia.tex
\section{Metodologia}
\label{sec:metodologia}

Nesta seção descrevemos as bases de dados utilizadas, as etapas de pré-processamento e construção de features, e o protocolo de avaliação adotado.

\subsection{Bases de Dados}

Trabalhamos com dez bases de dados de transações financeiras, cedidas no âmbito deste projeto como dados públicos. Elas estão organizadas em duas séries:

\begin{itemize}
    \item \textbf{Série PD}: bases \texttt{pd\_v0}, \texttt{pd\_v0.25}, \texttt{pd\_v0.5}, \texttt{pd\_v0.75} e \texttt{pd\_v1}.
    \item \textbf{Série GS}: bases \texttt{gs\_v0}, \texttt{gs\_v0.25}, \texttt{gs\_v0.5}, \texttt{gs\_v0.75} e \texttt{gs\_v1}.
\end{itemize}

O sufixo numérico de cada base indica seu \textbf{nível de viés controlado} ($v \in \{0{,}0;\; 0{,}25;\; 0{,}5;\; 0{,}75;\; 1{,}0\}$), onde $v = 0$ representa a ausência de viés introduzido e $v = 1$ corresponde ao nível máximo. Essa estrutura permite estudar sistematicamente como o grau de viés presente nos dados de treinamento e de teste afeta a equidade dos modelos.

O \textbf{atributo sensível} considerado é o ramo de atividade do estabelecimento comercial (\texttt{RAMO\_ATIVIDADE\_1}). Para manter a análise binária, filtramos apenas dois grupos: o \textbf{grupo privilegiado} (código 1) e o \textbf{grupo desprivilegiado} (código 4, recodificado para 0). A \textbf{variável alvo} é \texttt{I-d}, onde o valor 1 indica suspeita de fraude e o valor 0 indica transação legítima --- sendo este último o resultado considerado favorável nas métricas de equidade.

\subsection{Pré-processamento e Engenharia de Features}

Os dados são ordenados cronologicamente por titular antes de qualquer transformação, evitando vazamento de informação futura (\textit{data leakage}). A partir dos registros brutos, construímos dez features divididas em três grupos:

\begin{itemize}
    \item \textbf{Feature binária temporal:} \texttt{final\_de\_semana} indica se a transação ocorreu em sábado ou domingo, capturando um padrão comportamental distinto dos dias úteis.

    \item \textbf{Features de comportamento do titular} (calculadas com \textit{expanding windows} para evitar leakage):
    \begin{itemize}
        \item \texttt{pct\_saldo\_gasto}: proporção do valor da transação em relação ao saldo disponível;
        \item \texttt{zscore\_valor\_no\_titular}: z-score do valor da transação em relação ao histórico acumulado do próprio titular, sinalizando anomalias individuais;
        \item \texttt{n\_transacoes\_dia\_tit}: volume acumulado de transações do titular no mesmo dia;
        \item \texttt{is\_cnab220}: indicador binário de transações do tipo crédito (CNAB = 220);
        \item \texttt{prop\_cnab220\_tit}: proporção histórica de uso de crédito pelo titular, calculada sem vazamento (\textit{shift} de um período).
    \end{itemize}

    \item \textbf{Features de z-score por ramo de atividade} (neutralizam o efeito do grupo sensível ao medir anomalia relativa ao comportamento típico de cada ramo):
    \begin{itemize}
        \item \texttt{zscore\_valor\_no\_ramo}: anomalia do valor da transação dentro do ramo;
        \item \texttt{zscore\_n\_trans\_mes\_no\_ramo}: anomalia do volume mensal de transações dentro do ramo;
        \item \texttt{zscore\_reg\_no\_ramo}: anomalia da regularidade de uso do titular em relação ao seu ramo;
        \item \texttt{zscore\_valor\_total\_mes\_no\_ramo}: anomalia do volume financeiro mensal dentro do ramo.
    \end{itemize}
\end{itemize}

As features de z-score por ramo são particularmente relevantes para o problema de equidade: ao medir o comportamento do titular em relação à média do seu próprio grupo, elas reduzem a dependência do modelo em relação ao ramo de atividade como proxy discriminatório.

Antes do treinamento, apenas as features contínuas recebem padronização via \textit{StandardScaler}. As features binárias (\texttt{final\_de\_semana}, \texttt{is\_cnab220}) são mantidas em sua escala original. Valores ausentes são substituídos por zero.

\subsection{Protocolo de Avaliação}

\subsubsection{Balanceamento de classes}

Como é comum em detecção de fraudes, as bases apresentam forte desbalanceamento de classes. Para os experimentos que utilizam balanceamento, aplicamos subamostragem das classes majoritárias, limitando o conjunto de treinamento a no máximo 10.000 instâncias com até 5.000 por classe.

\subsubsection{Validação cruzada entre bases}

Adotamos um esquema de \textbf{avaliação cruzada entre bases}: cada base serve como conjunto de treinamento uma vez, e o modelo é avaliado em todas as demais. Isso nos permite observar como o viés do conjunto de treinamento influencia o desempenho e a equidade do modelo quando aplicado a ambientes com diferentes graus de viés. O conjunto de teste é limitado a 20.000 instâncias por base.

\subsubsection{Métricas de avaliação}

Os modelos são avaliados sob duas perspectivas. Do ponto de vista de \textbf{desempenho preditivo}:

\begin{itemize}
    \item \textbf{F1-Score}: média harmônica entre precisão e revocação, indicada para dados desbalanceados;
    \item \textbf{Precisão}: proporção de alertas de fraude que eram, de fato, fraudes;
    \item \textbf{Revocação}: proporção de fraudes reais corretamente identificadas.
\end{itemize}

Do ponto de vista de \textbf{equidade}, considerando a aprovação da transação (classe~0) como o resultado favorável:

\begin{itemize}
    \item \textbf{SPD} (\textit{Statistical Parity Difference}): diferença entre a taxa de aprovação do grupo privilegiado ($S=1$) e do grupo desprivilegiado ($S=0$). Valores próximos de zero indicam maior equidade:
    \[
    \text{SPD} = P(\hat{Y} = 0 \mid S = 1) - P(\hat{Y} = 0 \mid S = 0)
    \]

    \item \textbf{DI} (\textit{Disparate Impact}): razão entre as taxas de aprovação dos dois grupos. O critério dos ``quatro quintos'' considera valores abaixo de 0,8 como indicativos de discriminação sistemática:
    \[
    \text{DI} = \frac{P(\hat{Y} = 0 \mid S = 0)}{P(\hat{Y} = 0 \mid S = 1)}
    \]
\end{itemize}

Todos os experimentos utilizam semente aleatória fixa (\texttt{SEED = 42}) e são paralelizados com a biblioteca \texttt{joblib} para garantir reprodutibilidade e eficiência computacional.